\clearpage
\onecolumn
\section{Appendix A. Candidate Dataset and Simulation Register}
\newcommand{\sourceurl}[2]{\href{#1}{#2}}

\begin{center}
\begin{minipage}{.82\textwidth}
The resources below define a candidate development corpus, not a claim that
every named dataset has already been downloaded, licensed, harmonized, or
accepted into a released training mixture. Admission is versioned and
query-specific. Every snapshot receives a source card containing a persistent
identifier, version and checksums, license or data-use agreement, consent and
redistribution constraints, participant and session structure, inclusion
criteria, preprocessing lineage, native coordinate frames, clock and event
alignment, spatial and temporal support, observation or intervention model,
known missingness, demographic and acquisition coverage, and separate bias
and variance estimates. BIDS or NWB representations are retained where
available, but conversion does not erase the native format or provenance.

Dataset size does not determine gradient weight. A source is allowed to
update only the modules and ports it observes or defensibly constrains;
incompatible sources may retain separate parameters or be rejected by a
negative-transfer gate. Participant-level splits precede session and trial
splits, repeated measures remain grouped, and site or device leakage is
tested explicitly. Public archives are discovery and governance surfaces,
not homogeneous datasets. Controlled and clinical sources remain subject to
their own approvals. Synthetic sources are always labeled as
simulator-conditioned evidence and are evaluated against empirical data
before they influence an individualized posterior.
\end{minipage}
\end{center}

\begin{table*}[!h]
\centering
\footnotesize
\setlength{\tabcolsep}{3pt}
\renewcommand{\arraystretch}{1.08}
\begin{tabularx}{\textwidth}{@{}p{.14\textwidth}p{.21\textwidth}p{.24\textwidth}p{.15\textwidth}X@{}}
\toprule
\textbf{Resource} & \textbf{Native evidence and support} &
\textbf{Planned SC-WBD role} & \textbf{Access/governance} &
\textbf{Dominant bias, variance, and exclusion tests}\\
\midrule
Human Connectome Project (Young Adult; Development; Aging)
\citep{vanessen2013} &
Structural, diffusion, resting/task fMRI; selected MEG and behavioral
phenotypes; cross-sectional population geometry &
Population surface and tract priors, multimodal observation calibration,
task and lifespan conditioning &
ConnectomeDB; open and restricted tiers; release-specific terms &
Healthy-cohort selection, family structure, scanner/protocol concentration,
tractography ambiguity; split by family and release\\

UK Biobank Imaging &
Large-cohort structural, diffusion and functional MRI with health, genotype
and longitudinal subsets &
Population variation, anatomy--phenotype associations, aging and
disease-aware priors; not individual causal ground truth &
Application, data-use agreement and secure handling &
Volunteer and survivorship bias, population structure, incidental disease,
protocol drift; ancestry/site/age calibration required\\

Cambridge Centre for Ageing and Neuroscience (Cam-CAN)
\citep{taylor2017camcan} &
Adult-lifespan MRI, MEG, cognition, physiology and demographic measures &
Cross-modal clocks, age-conditioned priors, variance decomposition across
the adult lifespan &
Managed research access and citation conditions &
Cross-sectional aging is not longitudinal change; task and recruitment
effects; age-stratified residual checks\\

BigBrain \citep{amunts2013bigbrain} &
Microscopic three-dimensional histological reconstruction of one post-mortem
human brain &
High-resolution laminar and cytoarchitectonic geometry; scale bridges to MRI
surfaces &
Public atlas under stated use terms &
One brain, fixation and sectioning deformation, incomplete correspondence to
in-vivo physiology; never a population mean\\

Julich-Brain \citep{amunts2020julich} &
Probabilistic cytoarchitectonic maps assembled across post-mortem brains &
Uncertain areal boundaries, regional phenotype priors and atlas crosswalks &
Public atlas/data portal; versioned maps &
Sampling and histological boundary uncertainty, atlas warp and
inter-individual variation; retain probabilities rather than labels\\

Allen Human Brain Atlas \citep{hawrylycz2012} &
Spatial transcriptomic microarray/RNA evidence with MRI and histology from a
small donor series &
Regional molecular and receptor-related priors, gene-gradient hypotheses and
uncertainty maps &
Allen Institute terms and attribution &
Sparse sampling, donor/hemisphere imbalance, post-mortem effects and batch
correction; no voxelwise subject-specific expression claim\\

OASIS-3 and ADNI &
Longitudinal multimodal aging and neurodegeneration imaging with
cognitive/clinical measures &
Disease- and aging-aware validation after a healthy adult base model;
progression and scanner-shift stress tests &
Registration, data-use agreements and diagnosis-sensitive governance &
Clinical ascertainment, medication, missing-not-at-random visits, diagnostic
circularity and site drift; never pooled silently with healthy cohorts\\
\bottomrule
\end{tabularx}
\caption{\textbf{Candidate anatomical, molecular, and population-imaging
sources.} Each source constrains a declared scale and validity domain.
Histological and transcriptomic atlases inform population priors but do not
supply an individual living brain's microscopic state.}
\label{tab:appendix-anatomy}
\end{table*}

\clearpage
\begin{table*}[!h]
\centering
\footnotesize
\setlength{\tabcolsep}{3pt}
\renewcommand{\arraystretch}{1.08}
\begin{tabularx}{\textwidth}{@{}p{.14\textwidth}p{.21\textwidth}p{.24\textwidth}p{.15\textwidth}X@{}}
\toprule
\textbf{Resource} & \textbf{Native evidence and support} &
\textbf{Planned SC-WBD role} & \textbf{Access/governance} &
\textbf{Dominant bias, variance, and exclusion tests}\\
\midrule
HCP task/rest fMRI and multimodal subsets \citep{vanessen2013} &
Population task and rest fMRI aligned with anatomy, diffusion, behavior and
selected MEG &
Hemodynamic heads, task-conditioned regional trajectories and cross-modal
population baselines &
ConnectomeDB release snapshots &
Shared preprocessing, family leakage, task compliance and neurovascular
variability; retain acquisition and family IDs\\

Midnight Scan Club \citep{gordon2017msc} &
Dense repeated resting/task fMRI and anatomy in a small set of individuals &
Within-person functional topography, session variance and personalization
benchmarks &
Public research release &
Very small and nonrepresentative cohort; precision within these individuals
is not population generality\\

MyConnectome \citep{poldrack2015myconnectome} &
Dense longitudinal single-participant MRI, physiology, behavior and molecular
measurements over more than one year &
Slow-state, session-drift and person-specific adaptation tests &
Public project artifacts under source terms &
Single-subject and self-experiment selection, changing context and protocol;
used as a longitudinal case, not a universal subject\\

Natural Scenes Dataset \citep{allen2022nsd} and BOLD5000
\citep{chang2019bold5000} &
Natural-image fMRI with repeated or large stimulus sets and image identities &
Retinotopic and ventral-stream encoding, visual-memory interfaces,
fine-to-coarse observation mappings &
Public datasets with stimulus-license constraints &
Small cohorts, repeated-image adaptation, 7T/site specificity and stimulus
sampling; image identity may not be redistributed uniformly\\

StudyForrest \citep{hanke2014studyforrest} &
Naturalistic movie-centered fMRI with multimodal extensions, annotations and
behavior &
Long-horizon audiovisual dynamics, event boundaries and cross-study
naturalistic validation &
Public/OpenNeuro components; version-specific licenses &
One film and cultural/language context, annotation dependence, inter-run
alignment and scanner effects\\

Narratives \citep{nastase2021narratives} and Le Petit Prince fMRI
\citep{li2022lpp} &
Story-listening fMRI with stimulus timing; multilingual naturalistic speech
in the latter &
Auditory--language dynamics, temporal receptive fields and
language-generalization tests &
Public repositories; stimulus terms retained &
Story, language and comprehension heterogeneity; shared semantic timing does
not imply identical latent dynamics\\

Wakeman--Henson \citep{wakeman2015} and MOUS
\citep{schoffelen2019mous} &
Paired MRI, EEG/MEG and/or fMRI around face or language tasks &
Cross-method observation operators, clock alignment and source-localization
benchmarks &
Public datasets with repository-specific terms &
Task-restricted effects, modality sessions may not be simultaneous,
source-model dependence and small effects\\

OpenNeuro &
BIDS-indexed studies spanning MRI, EEG, MEG, iEEG and perturbation paradigms &
Versioned discovery layer and external-validation bank; individual studies
receive separate source cards &
Study-specific consent and licenses; snapshots pinned by dataset/version &
Archive membership is not harmonization; metadata quality, duplicate
subjects, preprocessing and redistribution vary by study\\
\bottomrule
\end{tabularx}
\caption{\textbf{Candidate functional, naturalistic, multimodal, and
longitudinal sources.} Dense individual studies test personalization; large
cohorts constrain population variation; naturalistic studies exercise longer
causal context. None supplies a complete latent trajectory.}
\label{tab:appendix-functional}
\end{table*}

\clearpage
\begin{table*}[!h]
\centering
\footnotesize
\setlength{\tabcolsep}{3pt}
\renewcommand{\arraystretch}{1.08}
\begin{tabularx}{\textwidth}{@{}p{.14\textwidth}p{.21\textwidth}p{.24\textwidth}p{.15\textwidth}X@{}}
\toprule
\textbf{Resource} & \textbf{Native evidence and support} &
\textbf{Planned SC-WBD role} & \textbf{Access/governance} &
\textbf{Dominant bias, variance, and exclusion tests}\\
\midrule
HCP and Cam-CAN MEG \citep{vanessen2013,taylor2017camcan} &
Task/rest MEG with subject anatomy and population phenotypes &
Millisecond latent dynamics, spectral/cross-frequency losses and MRI-informed
forward models &
Release- and project-specific access &
Head position, environmental noise, inverse non-identifiability and
cohort/protocol differences\\

Wakeman--Henson and MOUS \citep{wakeman2015,schoffelen2019mous} &
EEG/MEG with MRI and controlled face or language events &
Sensor-to-source calibration, modality-consistency losses and language/vision
motifs &
Public research releases &
Event-locked task narrowness, session separation, trial rejection and
analysis flexibility\\

PhysioNet EEG Motor Movement/Imagery &
Scalp EEG during executed and imagined motor tasks with event markers &
Motor-imagery BCI baselines, cortical motor readout and session-transfer
tests &
Public PhysioNet terms and attribution &
Coarse montage, cue and movement artifacts, variable imagery strategy; not a
direct motor-cortical source label\\

ZuCo \citep{hollenstein2018zuco} &
Simultaneous high-density EEG and eye tracking during natural reading &
Gaze-conditioned language slices, oculomotor interfaces and word/event
alignment &
Research access and corpus terms &
Small sample, fixation-related overlap, ocular contamination and
text-specific effects; gaze and EEG clocks retained\\

THINGS-EEG and THINGS-EEG2
\citep{grootswagers2022things,gifford2022things2} &
Large image/concept sets with rapid visual EEG repetitions &
Rapid visual dynamics, representation and recurrence tests,
stimulus-aligned pretraining &
Public dataset/stimulus terms &
Rapid-presentation and repetition effects, sensor-space mixing, image-set
semantics and subject imbalance\\

Temple University Hospital EEG Corpus \citep{obeid2016tuh} &
Large de-identified clinical EEG with reports and heterogeneous montages &
Artifact, abnormality and device-robustness pretraining; clinical
out-of-domain evaluation &
Registration/data-use agreement; clinical governance &
Care-path selection, label/report noise, medication, age and montage
heterogeneity; never treated as healthy resting physiology\\

RAM and other consented intracranial memory studies &
iEEG and stimulation events in clinical participants performing memory tasks &
Local temporal dynamics, hippocampal/cortical interfaces and perturbational
validation &
Controlled access; highly sensitive clinical data &
Epilepsy, electrode placement by clinical need, medication, sparse sampling
and stimulation carryover; no whole-brain coverage claim\\

DANDI and NEMAR/OpenNeuro archives &
NWB- and BIDS-centered electrophysiology repositories, including EEG, MEG and
iEEG studies &
Study discovery, format validation and held-out external tests &
Dataset-specific licenses, consent and versions &
Archive-level pooling prohibited; species, preparation, task, coordinate and
preprocessing semantics must match declared ports\\
\bottomrule
\end{tabularx}
\caption{\textbf{Candidate electrophysiology and BCI sources.} Sensor
coordinates, references, head models, event clocks, artifact decisions, and
inverse assumptions remain in the episode record. Intracranial recordings
offer local precision but not unbiased or whole-brain sampling.}
\label{tab:appendix-electrophysiology}
\end{table*}

\clearpage
\begin{table*}[!h]
\centering
\footnotesize
\setlength{\tabcolsep}{3pt}
\renewcommand{\arraystretch}{1.08}
\begin{tabularx}{\textwidth}{@{}p{.14\textwidth}p{.21\textwidth}p{.24\textwidth}p{.15\textwidth}X@{}}
\toprule
\textbf{Resource} & \textbf{Native evidence and support} &
\textbf{Planned SC-WBD role} & \textbf{Access/governance} &
\textbf{Dominant bias, variance, and exclusion tests}\\
\midrule
OpenNeuro ds004024 &
TMS--EEG with MRI, fMRI and diffusion data in a
paired-associative-stimulation/connectivity paradigm &
Electric-field registration, immediate response, propagation and multimodal
perturbation benchmarks &
Version-pinned OpenNeuro dataset &
Protocol-specific excitability, coil pose/field uncertainty, artifacts and
modest sample size; subject-level splits\\

OpenNeuro ds003037 (SUDMEX\_TMS) &
Longitudinal clinical MRI/TMS study in cocaine-use disorder &
Clinical longitudinal shift and response-prediction stress test after core
validation &
Version-pinned OpenNeuro dataset &
Diagnosis, treatment, attrition, medication/substance history and site
specificity; not a healthy-population prior\\

FieldTrip TMS--EEG tutorial data &
Motor-cortex TMS--EEG example recordings and analysis conventions &
Pipeline fixture, artifact-removal sensitivity and reproducible integration
test &
Tutorial terms; not a population corpus &
Pedagogical selection and preprocessing dependence; never used to estimate
clinical efficacy\\

tFUS mVEP speller dataset &
EEG during transcranial-ultrasound modulation in a visual evoked-potential
BCI paradigm, reported for 25 healthy participants &
Acoustic-write/EEG-read interface, BCI response and sham/condition modeling &
Figshare record 25583334; license and metadata pinned &
Device/protocol specificity, acoustic registration, EEG artifacts, blinding
and replication; effect direction validated independently\\

tFUS--ESI resting EEG dataset &
Resting EEG around transcranial focused ultrasound and transcranial
electrical stimulation &
State-dependent perturbation trajectory and cross-intervention comparison &
\href{https://doi.org/10.6084/m9.figshare.31029691}{Figshare record 31029691};
version pinned &
New-dataset maturity, device and montage effects, intervention order, focal
exposure uncertainty and multiple comparisons\\

AMASS \citep{mahmood2019amass}, Ego4D, and EPIC-KITCHENS &
Human motion capture and egocentric audiovisual action streams &
Embodied visual--motor priors, action/event segmentation and synthetic
sensorimotor boundary generation &
Dataset-specific noncommercial or research terms; media licenses retained &
No neural measurement, capture/actor/camera bias, action taxonomy and
cultural coverage; supervision stops at body/environment ports\\

WESAD, BioVid, Ninapro \citep{atzori2014ninapro}, Sleep-EDF and selected
PhysioNet physiology &
Wearable autonomic signals, pain/affect protocols, EMG hand actions, sleep
EEG/EOG/EMG, cardiac and respiratory streams &
Interoceptive, nociceptive, autonomic, peripheral motor and sleep-state
interfaces &
Source-specific licenses; pain/clinical data receive enhanced governance &
Construct labels, acted/elicited affect, peripheral-to-central ambiguity,
demographic coverage, missing sensors and protocol confounds\\

TRIBE v2 \citep{dascoli2026tribe} &
Population multimodal-to-fMRI encoder producing fsaverage5 cortical
predictions from video, audio or text &
Quarantined optional distillation factor at perceptual, language, or
hemodynamic interfaces; disabled by default and never a subject likelihood &
Official weights/code under CC BY-NC 4.0; model version pinned &
Not new data: perception--imagery, population--individual, teacher-domain,
rendered-proxy, hemodynamic, report and registration discrepancy; retain only
after matched empirical ablations\\

TVB \citep{sanz2013}, SimNIBS, ROAST, k-Wave, BabelBrain and known-source
electromagnetic simulations &
Neural-mass trajectories; electromagnetic and acoustic fields; synthetic
EEG/MEG with known latent sources &
Solver tests, intervention operators, rare regimes, identifiability and
recovery benchmarks &
Software licenses, versions, meshes and solver settings archived &
Model discrepancy and numerical error dominate; validate field and
neural-response components separately and never count synthetic replications
as participants\\

Embodied simulators and SuperCognition prospective study &
Physics-based humanoid sensory/motor/body trajectories; planned 23-adult,
89-session calibrated multimodal protocol &
Boundary pretraining followed by repeated individual inference, synchronized
read/write calibration and prospective validation &
Simulator licenses; human study pending ethics, consent, privacy review and
actual enrollment &
Simulation-to-human gap; planned counts are targets, not collected evidence;
preregister endpoints and retain device/pose/clock calibration\\
\bottomrule
\end{tabularx}
\caption{\textbf{Candidate perturbational, embodied, physiological,
prospective, and model-based sources.} Perturbation data receive privileged
causal status only for the intervention and outcomes actually randomized or
measured. Distillation models and physics engines shape priors or test
recovery; they neither add empirical participants nor validate biology.}
\label{tab:appendix-perturbation}
\end{table*}

\clearpage
\begin{table*}[!h]
\centering
\footnotesize
\setlength{\tabcolsep}{3pt}
\renewcommand{\arraystretch}{1.08}
\begin{tabularx}{\textwidth}{@{}p{.14\textwidth}p{.21\textwidth}p{.24\textwidth}p{.15\textwidth}X@{}}
\toprule
\textbf{Resource} & \textbf{Native evidence and support} &
\textbf{Planned SC-WBD role} & \textbf{Access/governance} &
\textbf{Dominant bias, variance, and exclusion tests}\\
\midrule
ABCD and HBCD
(\sourceurl{https://abcdstudy.org/}{ABCD};
\sourceurl{https://hbcdstudy.org/data-sharing/}{HBCD}) &
Longitudinal development, MRI/dMRI/fMRI, EEG in HBCD, biospecimens,
wearables, cognition, language, caregiver interaction and environment &
Age-conditioned anatomy and dynamics; developmental boundary conditions;
multimodal and longitudinal missingness models &
Controlled research access; child data, family structure and heightened
consent/privacy obligations &
Development is not adult stationarity; scanner/site, socioeconomic,
caregiver and attrition bias; never initialize an adult model without age
conditioning\\

Healthy Brain Network
(\sourceurl{https://fcon_1000.projects.nitrc.org/indi/cmi_healthy_brain_network/}{portal};
\sourceurl{https://openneuro.org/datasets/ds005505/}{EEG}) &
Pediatric multimodal MRI, EEG, eye tracking, actigraphy, genetics, voice,
video, cognition and clinical/behavioral phenotypes &
Cross-modal source cards, developmental BCI and naturalistic observation
ports, voice--face--gaze synchronization &
Open and managed components; pediatric and clinical sensitivity &
Community recruitment enriched for concerns, age and diagnosis
heterogeneity, device/version drift; participant and family grouping\\

NKI--Rockland Sample
(\sourceurl{https://fcon_1000.projects.nitrc.org/indi/enhanced/}{NKI-RS}) &
Lifespan MRI, diffusion, resting fMRI, phenotypes and selected repeated
measures &
Population variance, lifespan generalization and site-transfer tests &
INDI/NITRC access with release-specific terms &
Regional recruitment, protocol waves and uneven repeated measures; do not
treat cross-sectional age effects as trajectories\\

Individual Brain Charting
(\sourceurl{https://openneuro.org/datasets/ds000244/}{OpenNeuro ds000244}) &
Repeated high-resolution task fMRI across many cognitive paradigms in a fixed
cohort &
Task-to-region routing, within-person cognitive topography and cross-task
latent reuse &
Versioned OpenNeuro/EBRAINS releases &
Small cohort, task ontology and session practice effects; contrast labels are
not primitive mental-state labels\\

Courtois NeuroMod
(\sourceurl{https://www.cneuromod.ca/}{project};
\sourceurl{https://hub.datalad.org/naturalistic-imaging/cneuromod-opendata}{DataLad}) &
Deeply sampled individual fMRI with movies, games, images, tasks, behavior
and stimulus annotations &
Long-horizon individualized dynamics, context switching, memory and
naturalistic prediction heads &
Open components with dataset-specific stimulus licenses &
Very few extensively scanned participants, stimulus copyright and repeated
exposure; excellent within-person evidence but weak population coverage\\

TractoInferno
(\sourceurl{https://www.nature.com/articles/s41597-022-01833-1}{dataset paper})
and Fiber Data Hub
(\sourceurl{https://brain.labsolver.org/}{portal}) &
Multi-site diffusion MRI, bundles and large collections of processed
tractograms from HCP, ABCD, INDI and OpenNeuro &
Tract endpoint distributions, bundle-recognition priors, tractometry feature
heads and site-robust structural uncertainty &
Open or source-governed snapshots; preserve parent scan and processing recipe &
Algorithmic tractograms are derived evidence; false paths, bundle-definition
leakage and duplicate parent scans require lineage-aware splits\\

ISMRM 2015 Tractography Challenge
(\sourceurl{https://tractometer.org/ismrm2015/}{Tractometer};
\sourceurl{https://doi.org/10.5281/zenodo.572345}{data}) &
Brain-like diffusion phantom with known long-range bundles and challenge
tractograms &
Recovery tests for tractography/tractometry operators and synthetic
structural-connectivity uncertainty calibration &
Open challenge data and Zenodo versions &
Known simulated truth is valuable for algorithm error, not biological
prevalence; test domain shift to in-vivo data\\

Hansen neurotransmitter maps and neuromaps
(\sourceurl{https://github.com/netneurolab/hansen_receptors}{receptor atlas};
\sourceurl{https://netneurolab.github.io/neuromaps/}{neuromaps}) &
PET-derived group maps for receptors, transporters and other cortical
features, plus registered cross-atlas transformations &
Regional gain, neuromodulatory phenotype and prior-covariance hypotheses;
map-to-surface calibration &
Open code/maps under source terms; map-specific citations retained &
Group averages, tracer/selectivity, spatial autocorrelation and donor
differences; never infer a subject's receptor density from atlas value alone\\

EBRAINS Julich receptor-density data
(\sourceurl{https://julich-brain-atlas.de/atlas/neurotransmitter-receptors}{portal}) &
Post-mortem receptor-density measurements by area, layer and receptor class &
Laminar/regional parameter priors and model-class comparisons for
neuromodulatory operators &
EBRAINS knowledge-graph terms and versioned records &
Post-mortem preparation, small donor sets and parcellation correspondence;
retain measurement units and sampling locations\\

ENIGMA, BrainMap, NeuroSynth and NeuroQuery &
Consortium summary statistics and coordinate/text-derived meta-analytic maps
across disorders, tasks and cognitive terms &
External population regularizers, hypothesis generation and evaluation-only
maps for regional coverage &
Resource-specific terms; summary data do not erase cohort governance &
Publication and coordinate-reporting bias, term circularity and study
dependence; meta-analytic labels cannot supervise instantaneous latent state\\
\bottomrule
\end{tabularx}
\caption{\textbf{Additional developmental, dense-individual, tractometric,
connectomic, molecular, and meta-analytic sources.} Structural derivatives
retain their parent scan, algorithm and parameters. Atlas and meta-analytic
maps constrain priors or evaluations; they do not become subject-level
measurements by registration.}
\label{tab:appendix-connectome-expanded}
\end{table*}

\clearpage
\begin{table*}[!h]
\centering
\footnotesize
\setlength{\tabcolsep}{3pt}
\renewcommand{\arraystretch}{1.08}
\begin{tabularx}{\textwidth}{@{}p{.14\textwidth}p{.21\textwidth}p{.24\textwidth}p{.15\textwidth}X@{}}
\toprule
\textbf{Resource} & \textbf{Native evidence and support} &
\textbf{Planned SC-WBD role} & \textbf{Access/governance} &
\textbf{Dominant bias, variance, and exclusion tests}\\
\midrule
GazeBase
(\sourceurl{https://doi.org/10.6084/m9.figshare.12912257}{repository}) &
Longitudinal monocular eye movements across oculomotor tasks and repeated
rounds &
Oculomotor state, subject/session variability, identity leakage tests and
calibration-drift models &
Figshare dataset terms &
College-age concentration, one apparatus and task battery; biometric
identifiability requires sensitive governance\\

OpenEDS, OpenEDS2020 and TEyeD &
High-rate eye-camera sequences, pupil/iris/eyelid segmentation and gaze
targets, including head-mounted configurations &
Instrument observation model from eye image to gaze/pupil state; blink and
tracking-failure likelihoods &
Dataset-specific research licenses &
Camera geometry and illumination dominate; trains the sensor head, not visual
attention or cortical intent\\

COCO-Search18
(\sourceurl{https://sites.google.com/view/cocosearch/}{project}) &
Goal-directed gaze trajectories over natural scenes during 18 object-search
tasks &
Task-conditioned fixation/saccade policies, visual search and
world-centered salience interfaces &
Research release plus COCO image terms &
Small participant set, laboratory search and target-category restriction;
mouse-free gaze but no neural measurement\\

SALICON and MIT/Tuebingen Saliency Benchmark
(\sourceurl{https://salicon.net/}{SALICON};
\sourceurl{https://saliency.tuebingen.ai/}{benchmark}) &
Large free-viewing/saliency annotations, including mouse-contingent proxy
measurements and laboratory eye tracking &
Population visual-attention proposal prior and negative control for
task-free gaze prediction &
Image and benchmark licenses vary &
Crowdsourced mouse traces are not eye movements; central bias and image-set
semantics; never train subject-specific oculomotor dynamics as truth\\

ZuCo, OneStop, GECO, Provo and Dundee reading corpora &
Eye movements over natural text; ZuCo adds synchronized EEG; some corpora add
comprehension and bilingual conditions &
Gaze--word alignment, reading policy, lexical prediction, regressions and
language/oculomotor boundary states &
Corpus-specific access and text copyright &
Language, typography, comprehension, population and eye-tracker differences;
fixation duration is multiply determined\\

Ego4D gaze/social subsets and Project Aria
(\sourceurl{https://ego4d-data.org/}{Ego4D}) &
Egocentric video/audio with eye tracking in subsets, social interaction,
hands, activities and synchronized sensors &
Head--eye--world coordinate transforms, active vision, social orienting and
naturalistic audiovisual boundary trajectories &
License, privacy and face/voice restrictions; subset-specific metadata &
Wearable camera viewpoint, missing gaze, privacy-driven redaction and weak
neural linkage; preserve camera intrinsics and time bases\\

EPIC-KITCHENS and HD-EPIC
(\sourceurl{https://epic-kitchens.github.io/}{project}) &
Unscripted egocentric kitchen video, audio, narrations, actions, hands,
objects and 3-D/digital-twin annotations in HD-EPIC &
Long-horizon visuomotor episodes, object permanence, hand--object contact,
action hierarchies and digestive-input event proposals &
Research licenses; participant and household privacy; food/media metadata &
Kitchen/task domain, cultural and nutritional coverage, narration timing and
no direct neural signal\\

AMASS, Human3.6M, CMU MoCap and Motion-X &
Unified or source-specific body kinematics, meshes, actions and
text/audio-conditioned motion &
Biomechanical pose priors, efference/proprioception trajectories, motor
decoder pretraining and simulator boundary conditions &
Dataset-specific licenses; some noncommercial &
Marker-to-body fitting, actor/activity coverage and non-neural control
signals; enforce anthropometry and dynamics consistency\\

Ninapro
(\sourceurl{https://ninapro.hevs.ch/}{portal}) &
Surface EMG with hand kinematics, inertial and clinical variables across
intact participants and amputees &
Peripheral motor observation, intended-action interfaces, muscle synergy and
prosthetic-control benchmarks &
Public research datasets with source-specific terms &
Electrode shift, fatigue, anatomy, amputation and session variability; EMG
does not uniquely identify descending motor command\\

Touch and Go
(\sourceurl{https://touch-and-go.github.io/}{project}) &
Paired egocentric video and tactile-sensor recordings from human exploration
of real objects &
Visual--tactile association, contact-event prediction and haptic boundary
pretraining for embodied simulation &
Research dataset and media terms &
Instrumented touch differs from skin afferents; collector policy and object
coverage; supervision terminates at tactile boundary ports\\
\bottomrule
\end{tabularx}
\caption{\textbf{Candidate gaze, active-vision, movement, proprioceptive, and
tactile sources.} These corpora can train sensor, action and body-boundary
modules even without brain recordings. Their gradients must stop at the
boundary unless paired neural evidence supplies an identifiable continuation.}
\label{tab:appendix-sensory-motor}
\end{table*}

\clearpage
\begin{table*}[!h]
\centering
\footnotesize
\setlength{\tabcolsep}{3pt}
\renewcommand{\arraystretch}{1.08}
\begin{tabularx}{\textwidth}{@{}p{.14\textwidth}p{.21\textwidth}p{.24\textwidth}p{.15\textwidth}X@{}}
\toprule
\textbf{Resource} & \textbf{Native evidence and support} &
\textbf{Planned SC-WBD role} & \textbf{Access/governance} &
\textbf{Dominant bias, variance, and exclusion tests}\\
\midrule
Switchboard and Fisher English
(\sourceurl{https://catalog.ldc.upenn.edu/LDC97S62}{Switchboard}) &
Two-channel telephone conversations with transcripts, speakers, topics,
turns and dialogue-act annotations in derived releases &
Turn taking, interruption, repair, pragmatic transition and person-specific
speech/language boundary priors &
LDC licensing and redistribution restrictions &
Telephone band limit, prompted topics, era/dialect coverage and no visual or
physiological context; cannot supervise neural language state directly\\

AMI Meeting Corpus
(\sourceurl{https://groups.inf.ed.ac.uk/ami/corpus/}{portal}) &
Synchronized close/far-field audio, multi-camera video, slides, pens,
whiteboard, transcripts, dialogue acts, head movement and meeting roles &
Multi-party self/other tracking, speaker attention, conversational memory,
room acoustics and collaborative-action dynamics &
Public CC components plus registration and signal-specific terms &
Mostly staged design meetings, dropped-frame repairs, role and room
regularities; retain channel maps and synchronization corrections\\

IEMOCAP
(\sourceurl{https://sail.usc.edu/iemocap/}{portal}) &
Dyadic speech, video, facial/head/hand motion capture, transcripts and
categorical/dimensional affect annotations &
Audio--facial--gesture alignment, expressive action, prosody and
affect-readout pretraining &
USC research release; identity and voice sensitivity &
Actors, scripted/improvised elicitation, annotator disagreement and
stereotyped labels; train expression/readout, not latent feeling as ground
truth\\

RECOLA
(\sourceurl{https://recola.human-ist.ch/}{portal}) &
Synchronized French dyadic collaboration with audio, video, ECG, EDA and
continuous affect annotations &
Social interaction, autonomic--expressive coupling, delayed annotator
observation models and multimodal clock alignment &
Academic/nonprofit access conditions &
Small cultural/language context, remote task and annotation lag; valence and
arousal are observer variables, not sufficient affective state\\

MSP-Podcast, MELD and CMU-MOSEI &
Naturalistic or media-derived speech/video/text with emotion, sentiment,
speaker and dialogue context annotations &
Prosody, discourse-conditioned expression and robust multimodal readout heads;
hard negatives for generic context-conditioned language &
Resource-specific research terms and media copyright &
Podcast/television editing, label ontology, demographic imbalance and acted
or public persona; do not use labels as limbic targets\\

LibriSpeech, Common Voice, VoxCeleb and Bridge2AI Voice &
Read speech, crowdsourced speech, identity-rich interview audio and
clinically linked voice features/raw audio under controlled access &
Acoustic/phonetic front ends, speaker/session nuisance separation and
voice-to-respiratory/articulatory boundary models &
Open, community or controlled tiers; voice is identifiable biometric data &
Read versus spontaneous speech, internet/media noise, identity leakage,
health-label confounding and language imbalance\\

AudioSet, VGGSound, FSD50K and AudioCaps &
Large environmental audio or audiovisual event corpora with weak/strong
labels and captions &
Auditory scene decomposition, event-object proposals and source-separation
priors before neural fine-tuning &
Media links and licenses may change; source snapshots and hashes retained &
Web selection, label noise, clip context loss and copyright; event class is
not auditory-cortical representation\\

SpokenCOCO, HowTo100M and YouCook2 audio--language subsets &
Spoken image captions or instructional audiovisual speech aligned to objects,
steps and actions &
Audio--vision--language grounding, temporal event alignment and action
prediction at environmental ports &
Dataset and underlying media licenses &
Caption/task bias, narration describes rather than causes action, automatic
transcription error and instructional-domain concentration\\

MUSAN, DEMAND, AIR and BUT ReverbDB &
Speech/music/noise mixtures, ambient noise and measured/simulated room impulse
responses &
Microphone channel, reverberation, noise and source-separation calibration;
augmentation with known nuisance factors &
Research licenses; room/device metadata retained &
Acoustic nuisance corpora improve robustness but contain no cognitive labels;
simulation-to-room mismatch and convolution artifacts\\

Ego4D social interaction, AVA-AVD and VoxConverse &
Naturalistic audiovisual multi-speaker scenes with diarization, gaze/social
or active-speaker annotations &
Speaker persistence, face--voice binding, conversational scene memory and
self/other assignment &
Face/voice privacy, media and benchmark-specific terms &
Public/performed settings, annotation ambiguity, missing participants and
camera selection; test cross-corpus generalization\\

SuperCognition longitudinal conversation protocol &
Planned consented participant conversations, self-reports, tasks, voice,
video/gaze, physiology and synchronized EEG where feasible &
Individual lexical/prosodic dynamics, memory and preference updates, report
observation model and session-to-session person-state inference &
Prospective ethics, explicit voice/language reuse consent, withdrawal and
private/public release tiers &
Interviewer prompting, demand characteristics, privacy, rehearsal and model
feedback loops; preregister held-out situations and never scrape private
conversation\\
\bottomrule
\end{tabularx}
\caption{\textbf{Candidate speech, audio, conversation, and social-interaction
sources.} Unpaired corpora train environmental, linguistic, articulatory,
expressive and self--other boundary dynamics. They become evidence about brain
dynamics only through measured paired modalities or prospectively validated
transfer.}
\label{tab:appendix-conversation}
\end{table*}

\clearpage
\begin{table*}[!h]
\centering
\footnotesize
\setlength{\tabcolsep}{3pt}
\renewcommand{\arraystretch}{1.08}
\begin{tabularx}{\textwidth}{@{}p{.14\textwidth}p{.21\textwidth}p{.24\textwidth}p{.15\textwidth}X@{}}
\toprule
\textbf{Resource} & \textbf{Native evidence and support} &
\textbf{Planned SC-WBD role} & \textbf{Access/governance} &
\textbf{Dominant bias, variance, and exclusion tests}\\
\midrule
PhysioNet MIMIC-III/IV Waveforms
(\sourceurl{https://physionet.org/content/mimic4wdb/}{MIMIC-IV Waveforms}) &
ICU ECG, PPG, respiration, invasive/non-invasive pressure and numerics linked
to clinical context in subsets &
Cardiorespiratory boundary dynamics, artifact/device models and
pathophysiological stress tests &
Credentialed PhysioNet access, training and data-use agreement &
Critically ill single-system cohorts, interventions, missingness, device and
alarm artifacts; never a healthy autonomic prior without domain conditioning\\

BIDMC, CapnoBase, Fantasia and PPG benchmark collections &
ECG/PPG, respiration, capnography and reference intervals under resting or
clinical conditions &
Pulse transit, respiratory sinus arrhythmia, cardiorespiratory coupling and
wearable observation calibration &
Public or research-specific dataset terms &
Small cohorts and differing sensors/sampling; alignment and reference quality
vary; separate physiology from monitor preprocessing\\

Sleep-EDF, MASS, SHHS, MESA Sleep and ISRUC &
Overnight PSG: EEG, EOG, EMG, ECG, respiration, oxygenation, stages and
clinical/lifestyle variables in selected cohorts &
Sleep-state transitions, autonomic--cortical coupling, circadian context and
long-horizon missingness/noise models &
Open to controlled access; health and sleep data are sensitive &
Scoring-version and technologist disagreement, apnea/medication, montage and
cohort differences; stage labels are coarse observations\\

WESAD, DEAP and MAHNOB-HCI &
Wearable ECG, EDA, EMG, temperature and respiration or EEG/peripheral physiology
during stress and affective media protocols &
Autonomic observation heads, arousal/readiness proposals and
subject-normalization benchmarks &
Dataset-specific research licenses &
Elicited/acted states, small homogeneous samples, sensor placement and label
circularity; no direct limbic-state labels\\

UK Biobank accelerometry/biomarkers and NHANES &
Population blood/urine biomarkers, anthropometry, nutrition, activity,
wearables and health context; imaging linkage in UK Biobank subsets &
Slow systemic state, metabolic/endocrine covariates and population priors for
body-to-brain boundary conditions &
UKB application/DUA; NHANES public-use and restricted components &
Cross-sectional confounding, assay batch, self-report, selection and causal
ambiguity; slow covariates cannot supervise fast neural trajectories\\

Three-channel surface EGG dataset
(\sourceurl{https://zenodo.org/records/3878435}{Zenodo 3878435}) &
Raw and preprocessed cutaneous electrogastrograms from 20 healthy adults &
Gastric electrical observation head, preprocessing audit and fasting/fed or
rhythm feature calibration where metadata support it &
Zenodo version/license and participant metadata &
Small sample, three channels, electrode geometry, motion and abdominal volume
conduction; retain raw data and preprocessing parameters\\

Body-surface gastric mapping normative studies
(\sourceurl{https://pubmed.ncbi.nlm.nih.gov/36534985/}{normative reference}) &
High-resolution abdominal electrical maps, standardized meal response,
spectral/rhythm/amplitude reference intervals and symptom profiles &
Population priors and calibration coefficients for gastric slow-wave
frequency, rhythm stability, meal response and sensor noise &
Published aggregate coefficients; individual data only under study/device
governance &
Device/protocol specificity, BMI adjustment, meal composition and proprietary
processing; coefficients are calibration priors, not universal truths\\

CellML/Physiome model repository
(\sourceurl{https://www.cellml.org/}{CellML}) and gastrointestinal slow-wave
models &
Unit-annotated reusable cellular, electrophysiological, endocrine,
cardiovascular and gastrointestinal mathematical models &
Candidate transfer functions and priors for stomach, cardiac, respiratory,
endocrine and homeostatic modules; dimensional-consistency tests &
Model-specific licenses, versions and original validation papers &
Published mechanisms vary in species, preparation and identifiability; use as
competing model classes with discrepancy, never as data\\

Gastrointestinal cine-MRI and pan-alimentary MRI studies &
Dynamic stomach volume, contractions, emptying and bowel motility with meal
timing; paired EGG in some protocols &
Mechanical gastric-state observation heads, electrical--mechanical delay and
meal-conditioned body dynamics &
Often study-controlled rather than fully open; negotiate data and protocol
access &
Sequence/segmentation dependence, supine scanner context, sparse cohorts and
no direct enteric neural recording\\

GTEx, Human Protein Atlas, Human Cell Atlas and HuBMAP &
Tissue/cell expression, histology, proteomics and spatial molecular maps
across organs and donors &
Slow organ phenotype, receptor/transport and cell-class priors for peripheral
and autonomic interfaces &
Portal-specific consent, controlled genotypes and redistribution terms &
Post-mortem or surgical tissue, batch and donor effects, organ sampling and
healthy/disease mixture; never infer a living subject's expression directly\\

Human Microbiome Project and curated metabolomics resources &
Microbial community, functional, metabolite and host metadata across body
sites and cohorts &
Very slow gut/environment boundary covariates and hypothesis generation for
immune/metabolic state coupling &
Controlled individual-level data and resource-specific terms &
Diet, geography, medications, batch and compositionality dominate; no direct
neural causal edges without intervention or longitudinal mediation evidence\\
\bottomrule
\end{tabularx}
\caption{\textbf{Candidate autonomic, sleep, gastrointestinal, endocrine,
metabolic, and peripheral physiological sources.} These signals constrain the
body with which the brain is coupled. Aggregate clinical coefficients and
mechanistic repositories are calibration or model-class sources, not
participant observations.}
\label{tab:appendix-physiology-expanded}
\end{table*}

\clearpage
\begin{table*}[!h]
\centering
\footnotesize
\setlength{\tabcolsep}{3pt}
\renewcommand{\arraystretch}{1.06}
\begin{tabularx}{\textwidth}{@{}p{.14\textwidth}p{.21\textwidth}p{.24\textwidth}p{.15\textwidth}X@{}}
\toprule
\textbf{Resource} & \textbf{Native evidence and support} &
\textbf{Planned SC-WBD role} & \textbf{Access/governance} &
\textbf{Dominant bias, variance, and exclusion tests}\\
\midrule
Pyrfume and DREAM Olfaction Prediction Challenge
(\sourceurl{https://pyrfume.org/}{Pyrfume};
\sourceurl{https://www.synapse.org/Synapse:syn2811262}{DREAM}) &
Curated odorant chemistry, receptor/biology and psychophysical archives;
DREAM links 476 molecules to ratings from 49 people &
Chemical-to-olfactory-boundary priors, concentration and person variability,
odor-object representations and tests of learned olfactory geometry &
Archive-specific source licenses; DREAM/Synapse registration and challenge
terms &
Sparse chemical space, descriptor ontology, concentration, anosmia, culture
and small repeated cohorts; ratings supervise reports, not olfactory-cortical
state\\

BioVid Heat Pain and X-ITE Pain
(\sourceurl{https://www.nit.ovgu.de/nit/en/BioVid-p-1358.html}{BioVid}) &
Individually calibrated thermal/electrical nociceptive stimuli with facial
video, ECG, EDA, EMG and self-report in controlled protocols &
Dose-to-peripheral-response and expressive-readout operators; pain-report
likelihoods; hard tests for nociception, autonomic response and expression &
Application or research agreement; biometric video and health data &
Healthy volunteers, repeated laboratory pain, expectation, habituation and
sex/age/site effects; stimulus, report and experienced pain remain distinct\\

BP4D+ multimodal emotion corpus
(\sourceurl{https://cove.thecvf.com/datasets/523}{portal}) &
RGB/thermal video, 2-D/3-D landmarks, facial action units and peripheral
physiology from 140 adults during elicitation tasks &
Facial-muscle and autonomic observation heads, expression dynamics and
cross-channel timing; never a direct target for an affective latent &
Controlled academic access; identifiable face and physiology &
Elicitation/task regularities, selective frame coding, demographic coverage
and annotator ontology; validate actions and signals separately from labels\\

Sea Hero Quest
(\sourceurl{https://doi.org/10.1038/s41598-023-30937-w}{study}) &
Large-scale virtual wayfinding and path-integration trajectories with device
and optional demographic/context variables &
Population navigation-policy priors, path-integration errors, hippocampal
interface stress tests and age/culture-conditioned behavioral likelihoods &
Research approval and data-access statement; use only consented releases &
Self-selection, touchscreen/game literacy, device, country, missing
demographics and virtual-world policy; behavior does not localize its neural
implementation\\

ExtraSensory and PAMAP2
(\sourceurl{https://extrasensory.ucsd.edu/}{ExtraSensory};
\sourceurl{https://archive.ics.uci.edu/dataset/231/pamap2+physical+activity+monitoring}{PAMAP2}) &
Phone/watch context streams in everyday life; laboratory IMU, temperature
and heart-rate trajectories over 18 activities in PAMAP2 &
Body/environment/activity boundary states, wearable observation models,
multirate missingness and locomotor/metabolic context for longer episodes &
Public research releases and source terms &
ExtraSensory has 60 largely student/research users; PAMAP2 has nine subjects;
self-report, device placement and protocol effects prohibit population truth\\

StudentLife
(\sourceurl{https://www.cs.dartmouth.edu/~xia/publication/mobilehealth17/}{study}) &
Ten-week smartphone sensing, activity, sleep, location, conversation proxies,
self-report and academic context from a class cohort &
Slow contextual state, routine change, missing-not-at-random sensing and
longitudinal boundary episodes for person-adaptation experiments &
Public research data under study terms; mobility and mental-health sensitivity &
Forty-eight mostly computer-science students, one term and one institution;
correlations and phone-derived proxies cannot supervise causal mood dynamics\\

Open Humans and OpenAPS Data Commons
(\sourceurl{https://www.openhumans.org/create/}{Open Humans}) &
Participant-directed combinations of genomes, wearables, GPS, HealthKit,
glucose, insulin and other voluntarily shared longitudinal streams &
Source-card and consent infrastructure; rare dense-person time series;
glucose--meal--activity--sleep boundary dynamics and personalization tests &
Project- and member-specific permissions; consent is granular and revocable &
Extreme self-selection, heterogeneous preprocessing, missingness and
intervention feedback; never treat community donation as blanket reuse consent\\

Food-101 and USDA FoodData Central
(\sourceurl{https://data.vision.ee.ethz.ch/cvl/datasets_extra/food-101/}{Food-101};
\sourceurl{https://fdc.nal.usda.gov/}{FoodData Central}) &
Food images and labels; versioned composition, nutrient, portion and branded
food records rather than participant physiology &
Visual meal-object proposals and uncertain nutrient boundary conditions for
gastric, metabolic and endocrine simulations when an observed meal is paired &
Food-101 media terms; USDA data are public-domain/CC0 with version tracking &
Image label and cuisine bias, recipe/portion ambiguity and branded-product
drift; nutrient lookup is a noisy covariate, not ingested dose or absorption\\
\bottomrule
\end{tabularx}
\caption{\textbf{Additional chemical-sensory, nociceptive, expressive,
navigation, ecological-context, and nutrition sources.} These resources widen
the admissible boundary conditions and behavioral tests without pretending
that an odor rating, facial label, phone context, game trajectory or food name
is a measured neural latent. Each source receives an explicit gradient-stop
boundary and an observation model appropriate to its native support.}
\label{tab:appendix-additional-boundaries}
\end{table*}

\clearpage
\begin{table*}[!h]
\centering
\footnotesize
\setlength{\tabcolsep}{3pt}
\renewcommand{\arraystretch}{1.08}
\begin{tabularx}{\textwidth}{@{}p{.14\textwidth}p{.21\textwidth}p{.24\textwidth}p{.15\textwidth}X@{}}
\toprule
\textbf{Calibration / reference source} &
\textbf{Parameters or artifacts supplied} &
\textbf{Planned SC-WBD use} &
\textbf{Version/governance requirement} &
\textbf{Uncertainty and validation requirement}\\
\midrule
BIDS coordinate-system specification
(\sourceurl{https://bids-specification.readthedocs.io/en/stable/appendices/coordinate-systems.html}{coordinates}),
NWB and DICOM metadata &
Origins, orientations, units, fiducials, sensor/electrode frames, image
affines, channel and clock metadata &
Typed coordinate and acquisition manifests; import validation before any
spatial loss or cross-dataset alignment &
Pin specification/version and retain unmapped native fields &
Standards document metadata but do not verify it; compare landmarks,
determinants, units and round-trip transforms against raw headers\\

TemplateFlow, FreeSurfer, fsaverage, fsLR and standard atlases &
Versioned templates, surfaces, parcellations, transforms and density variants &
Explicit atlas/view projections and multiresolution surface crosswalks &
Record template identifier, resolution, software and transform chain &
Registration residual, topology mismatch and interpolation smoothing; atlas
coordinates never replace subject coordinates\\

IT'IS Tissue Properties Database
(\sourceurl{https://itis.swiss/virtual-population/tissue-properties/database}{portal}),
Gabriel/IFAC and NICT dielectric tables &
Frequency-dependent conductivity/permittivity, density, thermal, acoustic,
viscosity, relaxation and water-content parameters &
Priors for EEG/MEG lead fields, TMS/tES electric fields, tFUS propagation,
thermal exposure and MRI contrasts &
Archive database version, frequency, temperature, tissue definition and
citation; commercial terms checked &
Literature aggregation, ex-vivo/in-vivo and age variation; propagate
parameter distributions and compare alternate databases\\

SimNIBS head, coil and electrode examples and coil library
(\sourceurl{https://simnibs.github.io/simnibs/build/html/dataset.html}{datasets}) &
Segmented head meshes, standard electrode placement, TMS coil geometries and
finite-element reference outputs &
Unit/solver fixtures, pose-to-field operator tests and subject-head-model
pipeline initialization &
Pin SimNIBS, mesh, coil file, conductivities and solver tolerances &
Example outputs validate implementation consistency, not biological response;
compare mesh convergence and independent solvers\\

MNE-Python, FieldTrip and OpenMEEG reference datasets &
Forward/inverse examples, sensor geometries, spherical/BEM/FEM models,
filtering and event pipelines &
EEG/MEG lead-field, reference, source-localization and artifact regression
fixtures &
Version code, montage, head model, filters and rank decisions &
Inverse solution is non-unique and toolbox defaults differ; require
simulation recovery and cross-tool comparison\\

NIST/ADNI MRI phantoms, MRIQC and fMRIPrep reference outputs &
Geometry, signal uniformity, SNR, distortion, temporal stability, motion and
preprocessing quality metrics &
Scanner/session nuisance priors, QA rejection and uncertainty inflation for
MRI/fMRI observation heads &
Retain phantom serial/protocol, scanner software and pipeline/container hash &
Phantom stability does not calibrate neurovascular coupling; site thresholds
must be empirically related to downstream error\\

Breath-hold/hypercapnia CVR, arterial-spin-labeling and calibrated-fMRI data &
Vascular reactivity, perfusion, delay and oxygenation responses under known
gas or respiratory tasks &
Subject/session hemodynamic parameters and uncertainty for neural-to-BOLD
observation models &
Protocol, gas delivery, end-tidal traces, task compliance and physiological
monitoring retained &
Vascular and behavioral compliance effects; use to calibrate observation
heads, not neural state\\

Tractometer, TractSeg, pyAFQ, DIPY/RecoBundles and MRtrix fixtures &
Bundle definitions, streamline metrics, tract profiles, synthetic truth and
reference implementations &
Tractometry feature extraction, algorithm-disagreement ensembles and
connectome prior uncertainty &
Pin atlas/bundle definitions, seeding, thresholds, software and parent DWI &
Bundles are pipeline-dependent; compare algorithms and known-truth phantoms,
and prevent derivative duplication in splits\\

k-Wave, BabelBrain, acoustic phantoms and hydrophone measurements &
Acoustic propagation, transducer geometry, phase/steering, pressure and
thermal predictions or measurements &
tFUS drive-to-in-situ-field operator, solver discrepancy and device-specific
calibration &
Archive transducer, waveform, CT conversion, grid, boundary conditions and
hydrophone calibration &
Skull conversion, attenuation, coupling and cavitation uncertainty; validate
pressure fields before fitting neural effects\\

TMS--EEG artifact phantoms, coil probes and manufacturer waveform data &
Pulse waveform, coil current, induced-field measurements, amplifier recovery,
auditory/somatosensory artifact templates &
Separate physical dose, instrument saturation and peripheral co-stimulation
from candidate cortical response &
Device, coil, amplifier, ear masking and pose calibration versioned &
Vendor and laboratory specificity; sham rarely matches all sensory effects;
require artifact-injection recovery and negative controls\\

Gaze target grids, camera intrinsics and synchronization flashes &
Per-session eye-tracker mapping, pupil geometry, head pose, drift, latency and
clock offsets &
Calibrate eye-image/gaze observation operators and world/head/eye transform
chains before using fixation losses &
Raw calibration points and quality flags retained, not only vendor summary &
Calibration may be task- and gaze-range dependent; propagate spatially varying
error and detect recalibration breaks\\

Microphone calibration, room impulse responses and time-code/TTL tests &
Channel gain/phase, impulse response, array geometry, dropped frames, latency
and clock drift &
Acoustic observation model, diarization uncertainty and audiovisual event
alignment &
Retain hardware, sample-rate conversion, synchronization repairs and room
geometry &
Codec and automatic alignment can create apparent neural delays; validate with
physical sync events and known-source playbacks\\

Gastric electrode geometry, standardized meal and normative BSGM coefficients &
Electrode positions, channel transfer, slow-wave frequency/rhythm/amplitude,
meal timing and symptom-report alignment &
Gastroelectric observation model and slow body-state priors; initialize
individual calibration before longitudinal inference &
Record device, preprocessing, BMI/body geometry, meal composition and
reference-study version &
Device- and protocol-specific norms, volume conduction and proprietary
coefficients; measured raw signals override aggregate priors\\
\bottomrule
\end{tabularx}
\caption{\textbf{Candidate calibration, coordinate, physical-property, phantom,
and reference-implementation sources.} Calibration sources can train nuisance
parameters, observation/intervention operators and uncertainty. They do not
increase the number of biological observations and must not be mixed into
accuracy denominators as participant data.}
\label{tab:appendix-calibration-sources}
\end{table*}

\clearpage
\section{Appendix B. Machine-Readable Source Card and Gradient Contract}

A named dataset is not a training instruction. SC-WBD therefore ingests a
versioned \emph{source card} \(D_k\) that declares what a source measures, what
transformations were applied, and where its gradients are permitted to flow:

\[
D_k=\bigl(I_k,G_k,P_k,S_k,T_k,C_k,\mathcal O_k,\mathcal U_k,
M_k,B_k,V_k,\Delta_k,A_k,R_k\bigr).
\]

Here \(I_k\) is identity/version/provenance; \(G_k\) is governance;
\(P_k\) is the population and repeated-measures structure; \(S_k\) and \(T_k\)
are spatial and temporal support; \(C_k\) is the coordinate/calibration
manifest; \(\mathcal O_k\) and \(\mathcal U_k\) are observation and
intervention operators; \(M_k\) is missingness; \(B_k,V_k,\Delta_k\) are
systematic bias, sampling variance and model discrepancy; \(A_k\) is the
admissible parameter/module set; and \(R_k\) is the split and role policy. A
source that cannot populate a required field is not coerced into a common
raster. The field remains unknown, and the corresponding loss or gradient
path is disabled.

\begin{table*}[!h]
\centering
\footnotesize
\setlength{\tabcolsep}{4pt}
\renewcommand{\arraystretch}{1.10}
\begin{tabularx}{\textwidth}{@{}p{.15\textwidth}p{.23\textwidth}p{.28\textwidth}X@{}}
\toprule
\textbf{Source-card field} & \textbf{Required content} &
\textbf{Compiler consequence} & \textbf{Rejection or warning condition}\\
\midrule
Identity and lineage &
Persistent ID, release, file hashes, parent/derived relation, software and
container hashes &
Makes the exact training mixture reproducible and prevents a derivative from
crossing a parent-level split &
Unversioned mutable download, unknown parent, duplicate participant or
irreproducible preprocessing\\

Governance and consent &
License, DUA, purpose limits, consent scope, redistribution, withdrawal,
privacy and retention rules &
Controls storage, user visibility, allowed objectives and whether weights or
examples may be released &
Use outside consent/DUA, biometric data without safeguards, incompatible
license or unresolved withdrawal\\

Population and hierarchy &
Participant, family, site, device, session, run, trial, demographics and
selection mechanism &
Builds hierarchical effects and leakage-safe group splits &
Participant/family duplicated across folds, unknown selection or subgroup
too small for claimed calibration\\

Native signal and units &
Raw channels/voxels/events, units, reference, dynamic range, saturation and
quantization &
Instantiates typed ports and unit checks before resampling or normalization &
Unit ambiguity, undocumented reference, clipping or irreversible aggregation
that removes the target signal\\

Spatial support and point-spread &
Sensor/voxel/mesh/field support, resolution, kernel/lead field and sampling
geometry &
Attaches losses only at supported scales and creates projection uncertainty &
Nominal coordinate substituted for physical support or unjustified
fine-resolution label\\

Temporal support and clocks &
Sampling clock, exposure/integration window, filter group delay, event clock,
offset/drift and uncertainty &
Builds multirate frame maps, lag priors and synchronization factors &
Unknown clock relation, nonstationary dropped samples or alignment inferred
from the target effect without cross-validation\\

Coordinates and calibration &
Frame names, fiducials, transforms, calibration points, residuals, device
pose and raw metadata &
Creates a transform graph with propagated covariance rather than one baked
coordinate &
Missing units/origin/handedness, noninvertible transform or calibration
residual beyond query tolerance\\

Observation / intervention model &
Forward physics, preprocessing, dose, stimulus and which variables are
observed versus manipulated &
Separates read heads from write operators and permits causal losses only for
actual interventions &
Functional correlation mislabeled as intervention, unknown dose or
preprocessing that leaks labels\\

Missingness and censoring &
Planned/unplanned missing channels, artifact rejection, dropout, clinical
sampling and attrition &
Chooses marginalization, masks, inverse-probability/design factors or
missingness models &
Missing-as-zero labels, post-outcome selection or censoring correlated with
unmodeled state\\

Bias, variance and discrepancy &
Random noise, site/session variance, known systematic offsets, simulator or
teacher error and validity domain &
Determines uncertainty inflation, hierarchical weighting and whether a
source can update shared parameters &
Sample count used as reliability, zero teacher variance or pooled estimates
that erase site/subgroup bias\\

Gradient permission \(A_k\) &
Named modules, ports, scales and parameter groups the source may update;
frozen and evaluation-only objects &
Compiles an explicit gradient mask and audit log &
Loss reaches an unobserved module through an undeclared adapter or uses a
calibration source as biological supervision\\

Role and split policy \(R_k\) &
Prior, likelihood, boundary target, distillation, calibration, negative
control or evaluation-only; grouping and temporal holdout &
Selects the loss family and leakage barrier before model fitting &
Role changed after viewing test results, stimuli repeated across held-out
identity tests or synthetic replicas counted as independent subjects\\
\bottomrule
\end{tabularx}
\caption{\textbf{Minimum machine-readable source card.} The compiler converts
data governance and measurement semantics into executable gradient masks,
transform graphs, loss roles and leakage barriers. A source is rejected when
its semantics cannot support the intended claim.}
\label{tab:source-card}
\end{table*}

The permitted roles are deliberately non-equivalent. A \emph{likelihood}
factor contributes measured evidence about a latent variable through an
observation model. A \emph{boundary target} supervises a sensor, body or
environment port without licensing gradients through the unobserved brain.
A \emph{prior} shapes population parameters before individual evidence. A
\emph{distillation} factor transfers representational organization from a
teacher and carries teacher discrepancy. A \emph{calibration} factor estimates
coordinates, gains, tissue properties or nuisance parameters. A
\emph{negative control} should not predict the target under the claimed
mechanism. An \emph{evaluation-only} source never receives training gradients.
These roles are fixed in the run manifest before final evaluation.

\clearpage
\section{Appendix C. Coordinate, Clock, and Calibration Manifest}

Every sample and intervention is located in a graph of reference frames rather
than assigned one free-text coordinate. A manifest node declares a frame,
units, handedness, epoch and physical object; an edge declares a rigid,
affine, deformable, temporal or amplitude transform plus parameters,
calibration observations and uncertainty. A path is valid only if its
composed units and semantics agree. For a derived quantity \(z=f(x,c)\), where
\(c\) contains calibration coefficients, the ledger keeps first-order random
uncertainty and systematic offset separate:

\[
\Sigma_z \approx J_x\Sigma_xJ_x^\top+J_c\Sigma_cJ_c^\top
+J_x\Sigma_{xc}J_c^\top+J_c\Sigma_{cx}J_x^\top,
\qquad
b_z \approx J_xb_x+J_cb_c+\delta_f ,
\]

where \(\delta_f\) is model discrepancy. The cross terms are mandatory when
calibration error is shared across samples in a session; omitting them asserts
independence and can understate aggregate uncertainty. Monte Carlo or interval propagation
replaces this approximation for nonlinear registration, thresholded
tractography, acoustic focusing and other non-Gaussian maps.

\begin{table*}[!h]
\centering
\footnotesize
\setlength{\tabcolsep}{4pt}
\renewcommand{\arraystretch}{1.10}
\begin{tabularx}{\textwidth}{@{}p{.14\textwidth}p{.24\textwidth}p{.27\textwidth}X@{}}
\toprule
\textbf{Manifest layer} & \textbf{Stored quantities} &
\textbf{Examples} & \textbf{Required check before training/readout}\\
\midrule
Physical reference frame &
Object, origin, axes, handedness, units, validity interval &
Scanner RAS, voxel indices, subject surface RAS, fiducial head frame, tracker,
coil/transducer, cap, camera, eye and world frames &
Landmark and round-trip error; explicit distinction between continuous mm,
indices, normalized image coordinates and angular gaze\\

Deformation and projection &
Affine/nonlinear warp, surface registration, resampling kernel, Jacobian,
inverse availability and residual map &
Functional-to-anatomical, anatomical-to-MNI, fsnative-to-fsaverage/fsLR,
voxel-to-parcel, fine-to-coarse scale view &
No inverse assumed where absent; propagate local support and smoothing;
validate on held-out landmarks or correspondences\\

Sensor/device geometry &
Element centers/orientations, dimensions, camera intrinsics/extrinsics,
electrode/coil/transducer geometry and calibration date &
EEG contacts, MEG coils, TMS windings, tFUS elements, microphones,
eye cameras, gastric arrays and wearables &
Compare raw digitization/geometry to declared frame; reject swapped labels,
reflections and device files from another hardware revision\\

Clock graph &
Clock identity, sample rate, epoch, trigger path, offset, drift, jitter,
dropped samples and integration/filter delay &
Scanner volume/slice triggers, EEG amplifier, stimulus computer, eye tracker,
audio/video, neuronavigation and wearable clocks &
Physical synchronization event or independent cross-correlation target;
piecewise drift when one affine clock map is insufficient\\

Amplitude and unit calibration &
Gain, offset, frequency response, reference, impedance, quantization,
saturation and uncertainty &
EEG microvolts, MEG tesla, microphone pascals, EGG microvolts, PPG arbitrary
units, coil current, acoustic pressure and MRI intensity &
Dimensional check, calibration-range coverage and clipping detection; do not
standardize away physiologically meaningful amplitude without retaining map\\

Material/field parameters &
Tissue label, conductivity/permittivity, acoustic speed/attenuation, density,
thermal constants, anisotropy and source database &
EEG/TMS/tES FEM, tFUS propagation, thermal exposure and MRI relaxation &
Version/frequency/temperature match; sensitivity analysis across alternate
parameter tables and segmentation uncertainty\\

Stimulus and environment &
Display/audio calibration, luminance/color, sound level/impulse response,
camera pose, object/world geometry and event provenance &
Visual images, naturalistic video, speech, room conversation, VR, tactile
probe, standardized meal and embodied simulator &
Confirm delivered rather than requested stimulus; preserve media clock,
cropping/resampling, device response and environmental context\\

Intervention dose &
Waveform, timing, pose, field solution, target support, state, sham/control,
safety limits and operator uncertainty &
TMS pulses/trains, tFUS bursts, tES, sensory/cognitive task, neurofeedback and
drug/meal/body-state perturbations &
Dose and pose independently calibrated; randomized condition and sham
adequacy; separate field prediction from neural-response model\\

Calibration validity &
Calibration observations, fitting method, residuals, validity interval,
extrapolation distance and recalibration triggers &
Eye target grid, coil tracker, scanner phantom, hydrophone field,
body-surface gastric norms and microphone array geometry &
Do not reuse past validity interval silently; inflate uncertainty or require
recalibration when device, participant geometry or environment changes\\
\bottomrule
\end{tabularx}
\caption{\textbf{Coordinate, clock, and calibration manifest.} Cross-method
integration occurs through auditable transform paths with separate random
uncertainty, systematic bias and model discrepancy. The manifest prevents a
nominal atlas coordinate, sensor label or requested device setting from being
treated as the physical location or dose.}
\label{tab:calibration-manifest}
\end{table*}

\clearpage
\section{Appendix D. Heterogeneous Mixture, Leakage, and Evaluation Protocol}

For episode \(e\) from source \(k(e)\), SC-WBD composes only the loss factors
authorized by its source card:

\[
\mathcal L
=\sum_e w_e\!\left[
\lambda_{\rm like}\mathcal L^{(e)}_{\mathrm{likelihood}}
+\lambda_{\rm boundary}\mathcal L^{(e)}_{\mathrm{boundary}}
+\lambda_{\rm distill}\mathcal L^{(e)}_{\mathrm{distill}}
+\lambda_{\rm cal}\mathcal L^{(e)}_{\mathrm{cal}}
+\lambda_{\rm mechanism}\mathcal R^{(e)}_{\mathrm{mechanism}}
+\lambda_{\rm scale}\mathcal R^{(e)}_{\mathrm{scale}}
\right].
\]

The weight \(w_e\) is a learned or preregistered reliability term constrained
by effective sample size, participant/site hierarchy, estimated measurement
variance, systematic bias bounds, simulator/teacher discrepancy and overlap
with the target validity domain. It is not proportional to row count alone.
Losses are normalized within participant and source before cross-source
weighting so that a long recording or heavily oversampled simulator cannot
silently dominate. Gradient conflict is measured by module and source;
incompatible gradients trigger source-specific adapters, partial pooling,
freezing or rejection rather than forced consensus.

\begin{table*}[!h]
\centering
\footnotesize
\setlength{\tabcolsep}{4pt}
\renewcommand{\arraystretch}{1.10}
\begin{tabularx}{\textwidth}{@{}p{.15\textwidth}p{.22\textwidth}p{.29\textwidth}X@{}}
\toprule
\textbf{Failure mode or claim} & \textbf{Mandatory split/control} &
\textbf{Primary metric} & \textbf{Minimum interpretation}\\
\midrule
Participant or family leakage &
Group all sessions, derivatives, relatives and duplicate archive records
before splitting &
Held-out-person likelihood, calibration and retrieval/leakage audit &
Within-session prediction cannot support individual generalization\\

Stimulus memorization &
Hold out stimuli, semantic families and temporal continuations separately
from participant holdout &
Cross-stimulus neural/behavioral forecast and matched-feature baseline &
Image/audio recognition gain is not evidence for brain dynamics unless it
predicts new measured responses\\

Site/device shortcuts &
Leave-site/device/protocol-out evaluation; nuisance-only classifier and label
permutation within site &
Domain calibration, worst-site error and residual site predictability &
High pooled accuracy with poor new-device calibration is failure. On a
single-site corpus none of the three controls is constructible: the row is
\emph{not runnable}, which is a different report from \emph{not yet run}\\

Derived-data duplication &
Keep raw scan and every tractogram, parcellation, preprocessing derivative or
augmentation in one split &
Hash/lineage audit and performance after deduplication &
Different algorithms over one scan are not independent participants\\

Scale hallucination &
Withhold fine-scale evidence while retaining coarse data; compare uncertainty
and reconstruction to a coarse-only model &
Coverage and error at each native scale; high-frequency energy calibration &
Fine detail is valid only where source support or tested prior justifies it\\

Teacher/simulator domination &
No-teacher/no-simulator, generic-feature, shuffled, parameter-perturbed and
empirical-only ablations &
Measured held-out data likelihood and calibration, never teacher agreement
alone &
Distillation is retained only when it improves empirical prediction beyond
matched computation\\

Connectome prior value &
Randomized, distance-matched, dense, graph-only, local-only and soft-edge
controls at matched parameter/compute budgets &
Data efficiency, causal forecast, calibration and out-of-domain behavior &
Sparsity or plausibility alone is not evidence for the declared topology\\

Operator / mechanism claim &
Equal-capacity generic operator, alternate mechanism and learned-residual
controls; hold out discriminating perturbations &
Timing, direction, dose/state dependence and unique intervention forecast &
A mechanistic label is earned only by predictions a generic surrogate misses\\

Individualization claim &
Population, session-adapted, anatomy-only and longitudinal-person models; new
session and novel task/intervention holdouts &
Incremental log score, calibration, decision utility and drift &
Including a person's scan is not personalization unless it improves their
future predictions beyond population and session baselines\\

TMS/tFUS decision claim &
Prospective randomized or otherwise causally identified target/protocol
comparison with field, pose, state and sham records &
Directional response, dose--response, benefit/risk and decision regret &
Offline reconstruction supports target hypotheses, not wellness or treatment
efficacy\\

Language/person-model claim &
Hold out unfamiliar situations, future time windows and private facts;
compare generic LLM, language-history-only and multimodal person models &
Prospective choice/report/action calibration and counterfactual consistency &
Stylistic imitation or fact recall is not causal fidelity, consciousness or
personal identity\\

Dataset-family breadth &
Report performance by empirical, boundary-only, calibration, synthetic and
evaluation-only source roles; remove each family in turn &
Per-family contribution, negative transfer, subgroup worst case and
uncertainty coverage &
A longer source list is useful only when each family improves a specified
port or exposes a failure\\
\bottomrule
\end{tabularx}
\caption{\textbf{Minimum leakage controls, ablations, and interpretations for
heterogeneous training.} Every ambitious claim is paired with a split or
control capable of falsifying a shortcut explanation. The dataset register is
therefore an executable evaluation plan, not an assertion that more sources
automatically produce a better whole-brain model.}
\label{tab:mixture-evaluation}
\end{table*}

\textbf{``Not yet run'' and ``not runnable on this corpus'' are different
reports with different remedies.} Several controls in
Table~\ref{tab:mixture-evaluation} are constructible only on corpora that
contain the variation they test against, and a report that records both cases
as ``not passed'' conceals which is an engineering backlog item and which is an
acquisition requirement. The site/device row is the clearest case. It prescribes
three controls, and on a single-site, single-device corpus none of them exists:
leave-site-out requires a second site; a nuisance-only classifier requires site
labels that vary; and within-site permutation alone cannot falsify a site
shortcut, because the quantity it would have to permute is constant. No
splitting strategy repairs this, and no modeling change reaches it---the remedy
is a second site, and nothing else suffices. The same distinction applies to the
perturbation row of the claim-gate table in Section 0 when a corpus carries no
interventional structure: that is a data requirement, not a modeling one. Where
a control is unrunnable, the report must say so explicitly, name what would make
it runnable, and refrain from narrowing the associated claim on the strength of
a test that was never capable of running.

\twocolumn
